\chapter{Metodología}

La presente investigación adopta una metodología estructurada orientada a responder de manera sistemática las preguntas de investigación definidas para el estudio. Con este propósito, se diseñó un proceso compuesto por tres etapas principales: diseño de la revisión, búsqueda y selección de estudios, y análisis de evidencias.

Como se muestra en la Figura~\ref{fig:metodologia}, la primera etapa contempla la formulación de las preguntas de investigación, la definición de los criterios de selección y la elaboración del protocolo de revisión. La segunda etapa corresponde a la identificación, selección y filtrado de los estudios recuperados desde las bases de datos científicas, siguiendo un procedimiento documentado mediante el diagrama de flujo PRISMA, hasta obtener el corpus final de artículos incluidos en la revisión. Finalmente, la tercera etapa se enfoca en la extracción, clasificación y síntesis de la información contenida en los estudios seleccionados, utilizando matrices de análisis y categorías derivadas de las preguntas de investigación. Este proceso permite organizar la evidencia disponible, comparar los resultados reportados en la literatura e identificar brechas y oportunidades de investigación en el área de estudio.

% Diagrama Metodología
\begin{figure}[H]
\centering
\includegraphics[width=0.9\linewidth]{Resource/Metodología.png}
\caption{Diagrama de etapas de la metodología de investigación aplicada. Fuente: Elaboración propia.}
\label{fig:metodologia}
% Aquí va el diagrama PRISMA visual (TikZ)
\end{figure}
 
\section{Estrategia de búsqueda}
 
La investigación se estructura siguiendo la metodología PRISMA 2020
\cite{page2021prisma}, un marco de referencia ampliamente reconocido para la realización de revisiones sistemáticas de la literatura, que garantiza la transparencia, reproducibilidad y rigor científico del proceso de búsqueda y selección de estudios.
 
Para la selección de artículos científicos, se utilizaron dos bases de datos de amplio reconocimiento en el área de ciencias de la computación y ciberseguridad: \textbf{Scopus} y \textbf{Web of Science (WoS)}. La búsqueda se realizó mediante ecuaciones booleanas avanzadas que combinan términos relacionados con el análisis de logs de seguridad, el uso de modelos de lenguaje de gran escala (LLMs) y el razonamiento en ciberseguridad. Los conceptos clave empleados se presentan en la Tabla~\ref{tab:conceptos_clave}.

\begin{table}[H]
\centering
\renewcommand{\arraystretch}{1.2}
\caption{Conceptos clave de la búsqueda sistemática.}
\renewcommand{\tablename}{Table}
\setlength{\arrayrulewidth}{1pt}
\begin{tabularx}{0.9\textwidth}{@{} l X @{}}
\toprule
\hline
\textbf{Eje temático} & \textbf{Términos empleados} \\
\hline
\midrule
Análisis de logs & \textsl{log analysis, log parsing, log mining, event log analysis, system log, security log} \\
\hline
Modelos de lenguaje & \textsl{large language model, LLM, GPT, BERT, transformer model, foundation model} \\
\hline
Ciberseguridad & \textsl{cybersecurity, anomaly detection, threat detection, intrusion detection, SOC, SIEM} \\
\hline
Razonamiento & \textsl{causal reasoning, chain of thought, ReAct, multi-agent, explainable AI, XAI} \\
\bottomrule
\hline
\end{tabularx}
\label{tab:conceptos_clave}
\end{table}

\section{Preguntas de investigación}

Para guiar el proceso de revisión y asegurar coherencia con los objetivos del proyecto, se definieron cuatro preguntas de investigación con el propósito de abordar de manera integral los distintos componentes del problema de estudio (Tabla \ref{RQs}). La RQ1 se orienta a identificar las técnicas utilizadas para estructurar y analizar logs en sistemas complejos, así como sus principales limitaciones, proporcionando una base para comprender el estado del arte en términos metodológicos y las debilidades de los enfoques tradicionales. La RQ2 aborda el uso de modelos de lenguaje de gran escala en el análisis de logs, permitiendo caracterizar las tareas en las que estos modelos han sido aplicados y evaluar su aporte frente a técnicas previas.

Por su parte, la RQ3 se centra en los desafíos presentes en los centros de operaciones de seguridad, entregando información clave sobre las limitaciones operacionales del entorno real, tales como la sobrecarga de alertas, la dependencia de expertos y la complejidad del análisis de grandes volúmenes de datos. Finalmente, la RQ4 explora las capacidades de razonamiento incorporadas en los LLMs aplicados a ciberseguridad, tales como razonamiento causal, secuencial y explicativo, lo que permite evaluar el nivel de madurez de estas tecnologías y su potencial para abordar tareas complejas como la reconstrucción de incidentes. En conjunto, estas preguntas permiten no solo caracterizar el estado actual del conocimiento, sino también identificar brechas que justifican el desarrollo de nuevas soluciones basadas en agentes inteligentes.

\begin{table}[H]
\centering
\renewcommand{\arraystretch}{1.2}
\caption{Preguntas de investigación}
\renewcommand{\tablename}{Table}
\setlength{\arrayrulewidth}{1pt}
\begin{tabularx}{0.9\textwidth}{@{} l X @{}}
\toprule
\hline
\textbf{ID} & \textbf{Research question (RQ)} \\
\hline
\midrule
RQ1   & ¿Qué técnicas actuales se utilizan para estructurar y analizar logs en sistemas complejos y cuáles son las principales limitaciones? \\
RQ2   & ¿Cómo se están utilizando los modelos de LLMs para tareas de análisis de logs, tales como detección de anomalías, parsing y análisis de causa raíz? \\
RQ3   & ¿Cuáles son los principales desafíos que enfrentan los SOC para el análisis de logs de seguridad? \\
RQ4 & ¿Qué capacidades de razonamiento (causal, secuencial, explicativo) han sido incorporadas en modelos LLM aplicados a ciberseguridad y análisis de logs? \\
\bottomrule
\hline
\label{RQs}
\end{tabularx}
\end{table}

\section{Criterios de inclusión y exclusión}

Para asegurar la relevancia y calidad de los estudios seleccionados, se definieron los siguientes criterios:

\begin{table}[H]
\centering
\renewcommand{\arraystretch}{1.2}
\caption{Criterios de Inclusión}
\renewcommand{\tablename}{Table}
\setlength{\arrayrulewidth}{1pt}
\begin{tabularx}{0.9\textwidth}{@{} l X @{}}
\toprule
\hline
\textbf{ID} & \textbf{Criterios de Inclusión (CI)} \\
\hline
\midrule
CI-1   & Abordan el análisis de logs generados por sistemas computacionales o entornos de ciberseguridad. \\
CI-2   & Emplean técnicas de inteligencia artificial, incluyendo aprendizaje automático o modelos de lenguaje de gran escala (LLMs). \\
CI-3   & Aplicaciones en detección de anomalías, parsing o análisis de incidentes \\
CI-4   & Publicados en el periodo 2023 - 2026. \\
CI-5   & Publicados en fuentes académicas con revisión por pares.\\
CI-6   & Publicaciones en inglés. \\
CI-7   & Correspondientes a artículos de revista o proceedings de conferencias. \\
CI-8   & Publicaciones con un número mayor o igual a 3 citas. \\
\bottomrule
\hline
\end{tabularx}
\end{table}

% EXCLUSIÓN
\begin{table}[H]
\centering
\renewcommand{\arraystretch}{1.2}
\caption{Criterios de Exclusión}
\renewcommand{\tablename}{Table}
\setlength{\arrayrulewidth}{1pt}
\begin{tabularx}{0.9\textwidth}{@{} l X @{}}
\toprule
\hline
\textbf{ID} & \textbf{Criterios de Exclusión (CE)} \\
\hline
\midrule
CE-1   & No alineados con el análisis de logs en sistemas computacionales o ciberseguridad. \\
CE-2   & Ausencia de validación empírica mediante experimentación, métricas o evaluación cuantitativa. \\
CE-3   & Registros duplicados identificados en múltiples fuentes. \\
CE-4   & Publicaciones no académicas o sin proceso de revisión por pares. \\
CE-5   & Uso de modelos de lenguaje sin aplicación explícita al análisis de logs.\\
CE-6   & Sin disponibilidad de código o implementación abierta (non-open-source).\\
CE-7   & Documentos SLR.\\
\bottomrule
\hline
\end{tabularx}
\end{table}

\section{Fuentes de información y estrategia de búsqueda}

Como se indicó anteriormente, la búsqueda sistemática consideró dos fuentes de información científica: \textbf{Scopus} y \textbf{Web of Science (WoS)}. En cada una de ellas se aplicaron ecuaciones booleanas avanzadas, construidas a partir de términos asociados al análisis de logs, detección de anomalías, modelos de lenguaje y razonamiento en ciberseguridad.

Para garantizar una adecuada correspondencia entre los objetivos del estudio y la evidencia recuperada, se definió una ecuación de búsqueda específica para cada pregunta de investigación (RQ), evitando el uso de una única consulta general replicada en todas las bases de datos. Esta decisión metodológica permitió capturar con mayor precisión los conceptos particulares abordados por cada pregunta, mejorando la sensibilidad y especificidad del proceso de búsqueda. Las ecuaciones de búsqueda utilizadas y los resultados obtenidos por fuente de información se presentan en la Tabla~\ref{tab:busqueda_rqs}.


\begin{table}[H]
\centering
\renewcommand{\arraystretch}{1.1}
\renewcommand{\tablename}{Table}
\setlength{\arrayrulewidth}{0.5pt}
\begin{tabularx}{0.95\textwidth}{@{} l X r r r@{}}
\toprule
\hline
\textbf{RQ} & \textbf{Ecuación de búsqueda} & \textbf{Scopus} & \textbf{WoS}\\
\hline
\midrule
RQ1 &
\textsl{FIELD(("log analysis" OR "log analytics" OR "log mining" OR "event log analysis" OR 
"system log analysis" OR "software log analysis"))
AND
FIELD(("log parsing" OR "log structure*" OR "log template*" OR "log representation" OR 
"feature extraction" OR "log preprocess*" OR "pattern mining" OR "sequence mining" OR 
"anomaly detection" OR "unsupervised learning" OR "machine learning"))
AND
FIELD(("limitation*" OR "challenge*" OR "drawback*" OR "constraint*" OR 
"scalability issue*" OR "heterogeneous data" OR "data drift"))} & 61 & 40 \\[0.8em]
\hline
RQ2 &
\textsl{FIELD(("large language model*" OR LLM OR "transformer model*" OR GPT OR BERT OR 
"foundation model*"))
AND
FIELD(("log analysis" OR "log analytics" OR "log mining" OR "event log analysis" OR 
"system log*" OR "software log*"))
AND
FIELD(("anomaly detection" OR "log parsing" OR "template extraction" OR 
"root cause analysis" OR RCA OR "classification" OR "pattern recognition" OR 
"prediction" OR "fault diagnosis"))} & 75 & 41\\
\hline
RQ3 &
\textsl{FIELD(("security operations center" OR SOC OR "cybersecurity operation*" OR SIEM OR 
"security information and event management"))
AND
FIELD(("log analysis" OR "security log*" OR "event log*" OR "network log*" OR 
"threat detection"))
AND
FIELD(("challenge*" OR "limitation*" OR "alert fatigue" OR "data overload" OR 
"false positive*" OR "false negative*" OR "scalability" OR "complexity" OR 
"human factor" OR "analyst workload"))} & 43 & 27\\[0.8em]
\hline
RQ4 &
\textsl{FIELD(("large language model*" OR LLM OR GPT OR BERT OR "transformer model*"))
AND
FIELD(("reasoning" OR "causal reasoning" OR "sequential reasoning" OR 
"temporal reasoning" OR "explainable AI" OR XAI OR "interpretability" OR 
"inference" OR "decision making"))
AND
FIELD(("log analysis" OR "cybersecurity" OR "threat detection" OR 
"anomaly detection" OR "intrusion detection" OR "security log*" OR "event log*"))} & 44 & 26\\
\bottomrule
\end{tabularx}
\caption{Ecuaciones de búsqueda.}
\label{tab:busqueda_rqs}
\end{table}


\section{Proceso de Selección de Estudios}


El proceso de selección siguió las cuatro etapas del modelo PRISMA 2020: identificación, filtrado, elegibilidad e inclusión. Este procedimiento, sintetizado en la Figura~\ref{fig:prisma_flowchart}, permitió depurar progresivamente el conjunto inicial de registros mediante la aplicación de criterios de inclusión y exclusión orientados a asegurar la pertinencia temática, la calidad académica y la adecuación metodológica de los estudios seleccionados.

% Diagrama PRISMA
\begin{figure}[H]
\centering
\includegraphics[width=0.6\linewidth]{Resource/PRISMAv2.png}
\caption{Diagrama de flujo PRISMA 2020. Fuente: Elaboración propia.}
\label{fig:prisma_flowchart}
% Aquí va el diagrama PRISMA visual (TikZ)
\end{figure}

En la etapa de identificación, se recuperaron artículos desde dos fuentes de información científica: Scopus y Web of Science (WoS), mediante las ecuaciones de búsqueda previamente definidas. Como resultado, se obtuvieron 223 registros en Scopus y 134 en WoS, conformando un total inicial de 357 registros. Esta etapa permitió reunir un corpus amplio de publicaciones potencialmente relacionadas con el análisis de logs, los modelos de lenguaje de gran escala, la inteligencia artificial y la ciberseguridad.

En la etapa de filtrado, se aplicó inicialmente el criterio CE-3, correspondiente a la eliminación de registros duplicados identificados entre las fuentes de información. Como resultado, se removieron 148 documentos duplicados. Posteriormente, se aplicó el criterio CI-8, referido a la inclusión de publicaciones con un número mayor o igual a tres citas, lo que permitió excluir 124 registros adicionales mediante el filtro por citas. Tras esta depuración preliminar, el corpus se redujo a 85 documentos, los cuales fueron considerados para la etapa de elegibilidad.

En la etapa de elegibilidad, los 85 documentos restantes fueron evaluados en mayor detalle mediante la aplicación de los criterios de inclusión y exclusión definidos para la revisión. Se consideraron especialmente los criterios CI-1, CI-2 y CI-3, asociados a estudios que abordaran el análisis de logs en sistemas computacionales o entornos de ciberseguridad, que emplearan técnicas de inteligencia artificial o modelos de lenguaje de gran escala, y que se orientaran a tareas como detección de anomalías, parsing o análisis de incidentes. Asimismo, se verificó el cumplimiento de los criterios CI-4, CI-5, CI-6 y CI-7, relativos al periodo de publicación 2023--2026, revisión por pares, idioma inglés y tipo de publicación académica.

De manera complementaria, se aplicaron los criterios de exclusión CE-1, CE-2, CE-4, CE-5 y CE-6, descartando estudios no alineados con el análisis de logs en sistemas computacionales o ciberseguridad, trabajos sin validación empírica, publicaciones no académicas o sin revisión por pares, investigaciones centradas en modelos de lenguaje sin aplicación explícita al análisis de logs y documentos correspondientes a revisiones sistemáticas de literatura. Como resultado de esta evaluación, se excluyeron 33 documentos por no ajustarse a los criterios establecidos.

Finalmente, en la etapa de inclusión, se seleccionaron 52 estudios para conformar el corpus final de la revisión sistemática. Estos documentos constituyen la base del análisis posterior, al cumplir con los criterios definidos en términos de pertinencia temática, actualidad, respaldo académico y relación directa con la aplicación de inteligencia artificial y modelos de lenguaje al análisis de logs en contextos computacionales y de ciberseguridad.


Con el fin de garantizar la calidad y consistencia del conjunto de estudios analizados, se desarrolló un algoritmo de preprocesamiento orientado a la eliminación de registros duplicados y al filtrado de publicaciones según su impacto académico. Este procedimiento permitió depurar automáticamente los resultados obtenidos desde las bases de datos, conservando únicamente los registros más representativos y relevantes para la investigación. La aplicación de este algoritmo contribuyó a reducir la redundancia entre fuentes y a mejorar la calidad del corpus documental utilizado en las etapas posteriores de selección y análisis. El pseudocódigo correspondiente se presenta en el Algoritmo ~\ref{alg:duplicado_citas}.\\

\begin{algorithm}[H]
\setlength{\algomargin}{0.6em}
\caption{DOI duplicate removal and citation filtering}
\label{alg:duplicado_citas}

1. import pandas as pd\\
2. $rq\_por\_doi=df\_total.groupby('DOI')['RQ'].apply(unique)$\\
3. $fuente\_por\_doi=df\_total.groupby('DOI')['Fuente'].apply(unique)$\\
4. $df\_total\_ordenado=df\_total.sort\_values(['DOI','Citas'],ascending=[True,False])$\\
5. $df\_sin\_duplicados=df\_total\_ordenado.drop\_duplicates(subset= ['DOI'],keep='first')$\\
6. $df\_sin\_duplicados=merge(df\_sin\_duplicados,rq\_por\_doi)$\\
7. $df\_sin\_duplicados=merge(df\_sin\_duplicados,fuente\_por\_doi)$\\
8. $df\_filtrado=df\_sin\_duplicados[df\_sin\_duplicados['Citas']\geq MIN\_CITAS]$\\
9. $df\_filtrado.to\_csv('/content/filtered\_papers.csv',index=False)$

\end{algorithm}


\subsection{Evaluación del riesgo de sesgo}
No se aplicó una herramienta formal de evaluación del riesgo de sesgo, debido a que los estudios incluidos corresponden principalmente a investigaciones técnicas, propuestas metodológicas y evaluaciones experimentales sobre modelos de lenguaje de gran escala aplicados al análisis de logs y la ciberseguridad. La heterogeneidad de sus conjuntos de datos, métricas, entornos experimentales y objetivos dificulta el uso de instrumentos estandarizados para comparar directamente los resultados. En su lugar, la calidad metodológica de la revisión se resguardó mediante criterios de inclusión y exclusión rigurosos, la selección de publicaciones académicas revisadas por pares y un proceso de búsqueda, selección y síntesis transparente y reproducible, conforme a las directrices PRISMA 2020.


\section{Categorías de análisis}

Con el propósito de orientar de manera sistemática la revisión de la literatura planteada, se definió un sistema de categorías y subcategorías fundamentado en las preguntas de investigación formuladas. Esta estructura permite organizar la información recolectada de forma coherente con los objetivos del estudio, facilitando la identificación de enfoques, desafíos, capacidades técnicas y vacíos presentes en la literatura especializada.

La definición de las categorías se realiza a partir de las preguntas de investigación específicas, las cuales derivan de la pregunta principal: \textit{¿cómo un agente LLM puede analizar logs de seguridad y reconstruir incidentes?} A partir de esta pregunta central se establecen cuatro ejes de análisis, cada uno asociado a una categoría principal: 
\begin{itemize}
\begin{enumerate}
    \item Técnicas de análisis de logs (C1)
    \item Modelos LLM para logs (C2)
    \item Desafíos de SOC y SIEM (C3)
    \item Razonamiento en LLMs (C4)
\end{enumerate}
\end{itemize}


Se espera que estas categorías permitan abordar el problema desde una perspectiva técnica, operativa y analítica.

La primera categoría, técnicas de análisis de logs (C1), se vincula con la pregunta RQ1, orientada a identificar qué técnicas se utilizan para analizar logs en sistemas complejos. Dentro de esta categoría se consideran subcategorías como \textit{log parsing} (SC1.1) y detección de anomalías (SC1.2), ya que representan procesos fundamentales para estructurar, interpretar y detectar comportamientos relevantes en grandes volúmenes de registros generados por sistemas informáticos.

La segunda categoría, modelos LLM para logs (C2), se relaciona con la pregunta RQ2, enfocada en comprender cómo se emplean los modelos de lenguaje de gran escala en el análisis de logs de ciberseguridad. Esta categoría incluye subcategorías como \textit{fine-tuning} (SC2.1), \textit{prompt engineering} (SC2.2) y generación aumentada por recuperación, o RAG (SC2.3). Estas subcategorías permiten distinguir distintas estrategias de adaptación, consulta e integración de conocimiento externo en modelos LLM aplicados al análisis de eventos de seguridad.

La tercera categoría, desafíos del SOC y SIEM (C3), responde a la pregunta RQ3, centrada en identificar las principales dificultades que enfrentan los Centros de Operaciones de Seguridad en el análisis de logs. En este eje se consideran subcategorías como fatiga de alertas (SC3.1), correlación SIEM (SC3.2) y automatización SOC (SC3.3). Estas dimensiones permiten analizar las limitaciones operativas asociadas al manejo de grandes volúmenes de alertas, la integración de eventos provenientes de múltiples fuentes y la necesidad de automatizar procesos de detección y respuesta.

La cuarta categoría, razonamiento en LLMs (C4), se asocia con la pregunta RQ4, orientada a examinar qué capacidades de razonamiento presentan los LLMs aplicados a ciberseguridad. Esta categoría contempla subcategorías como razonamiento causal (SC4.1), \textit{chain-of-thought} (SC4.2) y explicabilidad mediante XAI (SC4.3). Su inclusión permite evaluar si estos modelos no solo identifican patrones en los logs, sino también si pueden establecer relaciones causales, construir explicaciones secuenciales y entregar resultados interpretables para apoyar la toma de decisiones en contextos de seguridad.

A partir de este sistema de categorías y subcategorías, los artículos seleccionados serán analizados de manera estructurada, identificando en cada estudio la información relevante asociada a los distintos ejes definidos. Este proceso permitirá clasificar y comparar los enfoques reportados en la literatura, así como reconocer tendencias, capacidades, limitaciones y desafíos vinculados al análisis de logs mediante modelos de lenguaje. De esta forma, el esquema categorial constituye el marco de referencia para la extracción y organización de la información, contribuyendo a la construcción sistemática del Estado del Arte y al análisis crítico de los antecedentes identificados.

Finalmente, la Figura~\ref{fig:categorias} presenta el esquema que relaciona la pregunta principal de investigación con las preguntas específicas, las categorías de análisis y sus respectivas subcategorías. Esta representación permite visualizar la lógica metodológica de la revisión, evidenciando cómo cada pregunta de investigación se vincula con una categoría analítica y cómo esta se descompone en dimensiones específicas que orientan el proceso de extracción, organización y análisis de la información.

\begin{figure}[H]
    \centering
    \includegraphics[width=1\linewidth]{Resource/DIAGRAMA_CATEGORIAS.drawio.png}
    \caption{Diagrama de categorías derivadas de las preguntas de investigación. Fuente: Elaboración propia}
    \label{fig:categorias}
\end{figure}

Las Tablas~\ref{tab:categorias_p1} y \ref{tab:categorias_p2} presentan las categorías y subcategorías relacionadas con las técnicas de análisis, modelos LLM, desafíos operacionales y capacidades de razonamiento empleadas en la literatura revisada.

% ── TABLA PARTE 1: C1 y C2 ──────────────────────────────────────────────
\begin{table}[H]
\centering
\renewcommand{\arraystretch}{1.4}
\caption{Categorías y subcategorías de análisis (Parte 1).}
\renewcommand{\tablename}{Table}
\setlength{\arrayrulewidth}{1pt}
\begin{tabularx}{0.9\textwidth}{@{} p{2.8cm} p{3.2cm} X @{}}
\toprule
\hline
\textbf{Categoría} & \textbf{Subcategoría} & \textbf{Definición} \\
\hline
\midrule

C1. Técnicas de análisis de logs
& SC1.1 Log parsing
& Proceso de transformar mensajes de log no estructurados en plantillas
  estructuradas y parámetros variables, habilitando el análisis
  automatizado posterior \parencite{ma2024llmparser}. \\
\cline{2-3}
& SC1.2 Anomaly detection
& Identificación de comportamientos que se desvían significativamente
  del patrón normal en secuencias de logs, indicando potenciales
  fallos o intrusiones \parencite{almodovar2024logfit}. \\
\hline

C2. Modelos LLM para logs
& SC2.1 Fine-tuning
& Proceso de adaptación de un LLM preentrenado a un dominio específico
  mediante entrenamiento adicional sobre datos del dominio objetivo,
  mejorando su rendimiento en tareas especializadas
  \parencite{karlsen2023llm4sec}. \\
\cline{2-3}
& SC2.2 Prompt engineering
& Diseño sistemático de instrucciones y contextos para guiar el
  comportamiento de un LLM hacia tareas específicas de análisis,
  sin necesidad de reentrenamiento \parencite{liu2024logprompt}. \\
\cline{2-3}
& SC2.3 RAG
& Técnica que combina recuperación de información relevante desde una
  base de conocimiento externa con la generación de respuestas del LLM,
  mejorando la precisión y reduciendo alucinaciones
  \cite{blefari2026cyberrag}. \\

\bottomrule
\hline
\end{tabularx}
\label{tab:categorias_p1}
\end{table}

% ── TABLA PARTE 2: C3 y C4 ──────────────────────────────────────────────
\begin{table}[H]
\centering
\renewcommand{\arraystretch}{1.4}
\caption{Categorías y subcategorías de análisis (Parte 2).}
\renewcommand{\tablename}{Table}
\setlength{\arrayrulewidth}{1pt}
\begin{tabularx}{0.9\textwidth}{@{} p{2.8cm} p{3.2cm} X @{}}
\toprule
\hline
\textbf{Categoría} & \textbf{Subcategoría} & \textbf{Definición} \\
\hline
\midrule
 
C3. Desafíos del SOC y SIEM
& SC3.1 Alert fatigue
& Fenómeno por el cual el exceso de alertas generadas por sistemas
  IDS/SIEM sobrepasa la capacidad analítica de los equipos de
  seguridad, aumentando el riesgo de ignorar incidentes reales
  \cite{palma2025llm}. \\
\cline{2-3}
& SC3.2 Correlación SIEM
& Proceso de relacionar eventos de seguridad provenientes de múltiples
  fuentes para identificar patrones de ataque multi-fase que no son
  detectables analizando eventos de forma aislada
  \cite{oduola2025wazuh}. \\
\cline{2-3}
& SC3.3 Automatización SOC
& Uso de tecnologías de inteligencia artificial para automatizar tareas
  repetitivas en el SOC, como el \textit{triage} de alertas, la
  clasificación de incidentes y la generación de reportes
  \cite{hassan2026elisar}. \\
\hline
 
C4. Razonamiento en LLMs
& SC4.1 Razonamiento causal
& Capacidad de un modelo para inferir relaciones de causa-efecto entre
  eventos de seguridad, permitiendo reconstruir la cadena de ataque y
  determinar el vector de entrada inicial
  \cite{liu2026ctithinker}. \\
\cline{2-3}
& SC4.2 Chain-of-thought
& Técnica de razonamiento paso a paso (\textit{CoT}) que permite a los
  LLMs descomponer problemas complejos en pasos intermedios, mejorando
  la precisión en tareas de análisis de seguridad
  \cite{hassan2026elisar}. \\
\cline{2-3}
& SC4.3 Explicabilidad (XAI)
& Capacidad de un sistema de inteligencia artificial para generar
  explicaciones comprensibles sobre sus decisiones, esencial para la
  trazabilidad y auditoría en contextos de ciberseguridad
  \cite{wang2026sentisec}. \\
 
\bottomrule
\hline
\end{tabularx}
\label{tab:categorias_p2}
\end{table}

